%%%%%%%%%%%%%%%%%%%%%%%%%%%%%%%%%%%%%%%%%%%%%%%%%%%%%%%%%%%%%%%%%%%%%%%%%%%%%%%%%%
\begin{frame}[fragile]\frametitle{}
\begin{center}
{\Large Risks \& Guardrails}
\end{center}
\end{frame}

%%%%%%%%%%%%%%%%%%%%%%%%%%%%%%%%%%%%%%%%%%%%%%%%%%%
\begin{frame}[fragile]\frametitle{Four Failure Modes Worth Knowing}
\begin{itemize}
\item \textbf{Hallucination:} confidently wrong answers, covered on the previous slide, always fact-check anything going to students or parents.
\item \textbf{Academic integrity:} students using AI to write essays undetected. Detection tools are unreliable; redesigning assessments (in-class writing, oral defense, process artifacts) works better than an arms race.
\item \textbf{Over-reliance:} treating the first draft as the final answer, for you or for students. AI output is a starting point, not a verdict.
\item \textbf{Bias:} models trained on internet-scale text inherit internet-scale biases. Review anything touching grading, feedback tone, or student evaluation with that in mind.
\end{itemize}
\end{frame}

%%%%%%%%%%%%%%%%%%%%%%%%%%%%%%%%%%%%%%%%%%%%%%%%%%%
\begin{frame}[fragile]\frametitle{Student Data: The One Non-Negotiable}
\begin{itemize}
\item \textbf{Never paste} student names, grades, health information, or any personally identifiable data into a public, non-institutional AI tool.
\item If your school has an approved, FERPA-compliant (or local equivalent) tool, use that for anything involving real student records.
\item For everything else, anonymize first: ``a Grade 7 student who struggles with fractions'' works just as well as a real name for drafting feedback.
\end{itemize}
\begin{alertblock}{Rule of Thumb}
If you would not write it on a sticky note left on a public desk, do not paste it into a public AI chat window.
\end{alertblock}
\end{frame}

%%%%%%%%%%%%%%%%%%%%%%%%%%%%%%%%%%%%%%%%%%%%%%%%%%%
\begin{frame}[fragile]\frametitle{Quick Check: Spot the Problem}
\begin{itemize}
\item A colleague pastes this into a public AI chatbot to get feedback drafted faster:
\end{itemize}
\begin{lstlisting}[language=HTML]
Write encouraging feedback for Priya Sharma, roll no. 23,
who has been diagnosed with ADHD and is struggling with
her Grade 8 persuasive essay. Her essay: [pastes full text]
\end{lstlisting}
\vspace{15pt}
\begin{itemize}
\item What is wrong here, and what is the one-line fix?
\end{itemize}
\end{frame}

%%%%%%%%%%%%%%%%%%%%%%%%%%%%%%%%%%%%%%%%%%%%%%%%%%%
\begin{frame}[fragile]\frametitle{Quick Check: Answer}
\begin{itemize}
\item \textbf{The problem:} full name, roll number, and a health diagnosis, all personally identifiable, all sent to a public tool with no institutional data agreement.
\item \textbf{The fix:} drop the name, roll number, and diagnosis. Keep only what the task actually needs: ``a Grade 8 student who needs extra encouragement and structure'' plus the essay text.
\item Same quality of feedback, none of the exposure.
\end{itemize}
\end{frame}
